\documentclass[a4paper]{article}
\usepackage{MMVII}

\newcommand{\TheCmd}{ReportTieP}

\begin{document}

\CmdTitle
%\maketitle

\tableofcontents


 %=================================================================================

\SecDescr

This command generate report on quality of orientation relative to a set of tie points,
by default it generates these analysis by image using multiple tie-points, but it can optionnally generate
result for specific pairs.  It generate result similar to  \FileSty{ReportGCP}, see also
this command for detailled interpretation of files.

It generate three kind of information :

\begin{itemize}
  \item statisical information in \FileSty{csv} format , for indvidual images and camera, and possibly for pairs;


  \item $2D$ visual information, in  \FileSty{jpg/tiff} format of the residuals, that can be exported as 
       dense images of field of vectors;  

  \item $3D$ visual information, in  \FileSty{ply} format to see the global spatialisation of the residuals. 
\end{itemize}

This variety of export is required because the potential problem in orientation can have many different origin :
problem of internal calibration, problem due to one or several images, problem due to a specific pair of image , problem due
to some part of the 3d scene....
Generally, all these export are not generated simultaneously, the user can generate a first report without
optionnal parameter and then focuss and specific reports.


 %=================================================================================
\section{Run with default parameters}


An example of run With the three defautt parameters example on Ramses dataset:

\begin{verbatim}
MMVII ReportTieP IMG_03.*JPG V1Dense Adjust 
\end{verbatim}


It will generate statistic of residual of residual in folder \FileSty{MMVII-PhgrProj/Reports/ReportTieP/Adjust\_V1Dense}, 
it will contain :

\begin{itemize}
   \item  a file  \FileSty{Glob.csv} with statistic on all images;
   \item  a file  \FileSty{ByCam.csv} with statistic on each individual camera ; in the current case 
         where there is only one camera these two files will be identicall; 
   \item  a file  \FileSty{ByImage.csv} with statistic on each individual image.
\end{itemize}


The statisic contains residual for given percentile that can be modified with otional parameter  \FileSty{Perc}.
It will also generate a dense image for the used cameras in the folder  \FileSty{MMVII-PhgrProj/VISU/ReportTieP/Adjust\_V1Dens}
if the optional paramete  \FileSty{DoImResCalib} is set to  \FileSty{true} (its default value).


 %=================================================================================
\section{Optional parameters  for  per/image visualisation}

If we are interested by specific pair of images, we must use tie-point that comes
from initial pairs and not merge tie points. So : 

\begin{itemize}

   \item  statistic  on residual for all pairs of images will be computed, they
          will be saved in the file   \FileSty{ByPairs.csv} of the  report folder;
          they can be modified with otional parameter  \FileSty{Perc};

   \item  parameter  \FileSty{PatImFV} , is a pattern that is used as filter for
          selecting the image for which we generate the residual as a field
          of $2D$ vector superposed to each image.

   \item  parameter  \FileSty{PatImRes} , is a pattern that is used as filter for
          selecting pairs of image for which we generate dense visualisation;
\end{itemize}

An example of use with Ramses dataset :

\begin{verbatim}
MMVII ReportTieP .*JPG V1Dense Adjust PatImRes="IMG_035.*.JPG"  PatImFV="IMG_035[3-6].JPG"
\end{verbatim}

Will generate $10$ images of dense residual and $4$ images of vector field. The vector field
images is controled by parameter \FileSty{ParamFV} , see {\tt ReportGCP}.


 %=================================================================================
\section{Optional parameters  for  per/pair statistic and visualisation}

If we are interested by specific pair of images, we must use tie-point that comes
from initial pairs and not merged tie points. So : 

\begin{itemize}
    \item we must use parameter  \FileSty{InTieP} to activate computation by pairs and specify the folder containing the 
          tie points of pairs pairs;

    \item static for all pairs will be generated in file   \FileSty{ByPairs.csv}, using 
          parameter  \FileSty{Perc};

    \item the optional parameter \FileSty{PatImRes} specify the images for which we generate 
          dense visualization;

    \item the optional parameter \FileSty{PairPatFV} specify the images for which we generate 
          the pair for  which we generate visualisation of residuals as vector's field;

\end{itemize}

An example of use with Ramses dataset :

\begin{verbatim}
MMVII ReportTieP IMG_03.*JPG V1Dense Adjust InTieP=V1Dense  PairPatRes=IMG_035[0-3].JPG  PairPatFV=IMG_035[1-2].JPG
\end{verbatim}

will generate  file \FileSty{ByPairs.csv} with $213$ lines corresponding to all pairs having non empty set of tie points,
one image vector fields for \FileSty{(IMG\_0351.JPG,IMG\_0352.JPG)} and $6$ image dense
for pairs of  $\{50-51,50-52,\dots,52-53\}$.

 %=================================================================================

\section{Optional parameters  for  $3d$ visualisation}

The $3d$ visualisation cab be used to analyse potentiel correlation between localisation 
and residual. For this the parameter optional  \FileSty{ParamPly} can used :

\begin{itemize}
    \item  if it used, a file \FileSty{RTP.ply} will be generated ;

    \item  the $3d$ point will be computed by bundles intersection using ties point and orientations;

    \item the color of each point will depend of the residual .
           
\end{itemize}



The parameter  \FileSty{ParamPly} contains $2$ value, a boolean indicating if generate
a coloured of a gray file, and an exponent $E$.
For computing the colour :

\begin{itemize}
   \item  first we compute a global histogramm $H$ of residuals;
   \item  then for each point, the residual $r$ is converted to rank  $Rk$ with $Rk \in [0,1]$;
   \item  then $Rk$ is converted to a quality $Q$ using the formula $Q=1-Rk^E$ (some kind of \emph{gama-correction});
   \item  in gray mode, the colour is $Q$  (point with high residual are black)
   \item  in colour mode, $Q$ is used to compute an  abscissa,  on the path $Red-Orange-Green$ 
         (point with high residual are red).
\end{itemize}


A consequence, of this way of proceeding, is that due to histogramm pre-processing, if we multiply all the
residual by $2$ or $10$ or \dots , we will obtain exactly the same file.  This visualization is not
made for  quantitative evaluation, but for qualitative localisation of possible area with problem.






\end{document}

